% =====================================================================================================
\section{The posterior selects the right model}\label{sec:selection}

This section shows that the posterior puts almost all its mass on the true model, at the rate $\rn$, for fixed and for
random Gaussian designs. On every design no rate faster than $n^{-1/2}$ is possible, and on two families of designs the
rate $\rn$ itself is attained. When some true predictors are missing from the data, the posterior concentrates instead
on the predictors used by the best linear predictor built from the available ones.

\subsection{Fixed designs}\label{sec:fixed}

Our first result answers the basic question: with the design fixed, how much posterior mass lies outside the true
model?

\inwords{If the true model is small, the coefficients are bounded, small groups of predictors are not collinear, and every
true signal is at least a constant times $\rn$, then, on average over the noise, the posterior mass outside the true model
is at most a constant times $\rn$. The constant is the same for every design that meets the conditions. As a consequence,
the highest posterior model and the median probability model both equal the true model with probability tending to one,
and every inclusion probability is within order $\rn$ of $1$ for a true predictor and of $0$ for a null one.}

\begin{assumption}{SP}{Sparsity}
The true model size $\sz\ge1$ is fixed.
\rulesout{true models whose size grows with $n$ (Theorem~\ref{thm:sharp} shows what is possible then), and the empty
true model.}
\end{assumption}

\begin{assumption}{BC}{Bounded coefficients}
$\max_{j\in\set{p}}\abs{\theta^*_j}\le\bar\theta$ for a fixed $\bar\theta<\infty$.
\rulesout{true coefficients that grow with $n$. The signal energy $\norm{\mu}^2$ then stays of order $n$, like the noise
energy.}
\end{assumption}

\begin{assumption}{EF}{Eigenvalue floor on sets of bounded size}
$\lmin(\Rhat_S)\ge\varphi_0^2$ for every model $S$ with $\abs{S}\le2(A+1)\sz+2$, for a fixed $\varphi_0\in(0,1]$.
\rulesout{strong collinearity within small sets of predictors, which would let a wrong small model mimic the true one.}
\end{assumption}

\begin{assumption}{NC}{No near-exact collinearity among sets of logarithmic size}
$\lmin(G_S)\ge g_0$ for every model $S$ with $\abs{S}\le3(A+1)(\sz+1)\log n$, for a fixed $g_0>0$.
\rulesout{exact, or nearly exact, linear dependence among $O(\log n)$ predictors; duplicated predictors, for instance.}
\end{assumption}

On the correlation scale \as{NC} reads $\lmin(\Rhat_S)\ge g_0/(n-1)$: a floor that vanishes as $n$ grows. We call \as{NC}
the \emph{vanishing floor}.

\begin{assumption}{BM}{Beta-min}
Either (a) $\tmin\ge C_\beta\,\sigs\rn$ with a fixed $C_\beta>0$ such that $C_\beta^2\varphi_0^2>c_0(\sz,A)$, where
$c_0(\sz,A)$ is an explicit constant that depends only on $(\sz,A)$; or (b) $\tmin\sqrt{n/\log n}\to\infty$.
\rulesout{true signals too small to be told apart from noise at the scale $\rn$.}
\end{assumption}

\begin{theorem}[Posterior concentration, fixed design]\label{thm:post-fixed}
Let the design be fixed and assume \as{GL}, \as{DIM}, \as{SP}, \as{BC}, \as{EF}, \as{NC} and \as{BM}. There are $n_0$
and $C_1$, depending only on $(\sz,A,\sigs^2,\bar\theta,\varphi_0,g_0,C_\beta)$, such that for every $n\ge n_0$ and every
design and $\tstar$ that satisfy these assumptions at that $n$,
\begin{equation}\label{eq:post-exp}
  \E\big[1-\pis{\Sz}\big]\ \le\ C_1\sqrt{\frac{\log n}{n}} .
\end{equation}
Consequently, along every sequence that satisfies the assumptions,
\begin{gather}
  \P\big(\HPM=\MPM=\Sz\big)\ \ge\ 1-4C_1\sqrt{\frac{\log n}{n}}\ \to1,\label{eq:post-hpm}\\
  \max_{j\in\set{p}}\big|\PIP_j-\ind{j\in\Sz}\big|\ \le\ 1-\pis{\Sz}\ =\ \OP\Big(\sqrt{\frac{\log n}{n}}\Big).\label{eq:post-op}
\end{gather}
The expectation is over the noise, with the design fixed. Under form (b) of \as{BM}, $n_0$ also depends on how fast
$\tmin\sqrt{n/\log n}$ diverges.
\end{theorem}

\status{The in-probability conclusions (the limit in~\eqref{eq:post-hpm}, and~\eqref{eq:post-op}): \textbf{proved and independently verified} (two independent verification passes: one by a reviewer of the same model
family as the author, one by a reviewer of a different model family; all inputs of the result were discharged on
2026-11-16). The
expectation bound~\eqref{eq:post-exp}, and the explicit bound $1-4C_1\rn$ in~\eqref{eq:post-hpm} that follows from it:
\textbf{proved and checked, conditional on two supporting results} (a proposition giving a lower bound on the posterior
miss, and a lemma bounding the tail of the posterior model size) \textbf{that are proved but not yet independently
checked}. \bridgepending{} Verbatim registry strings: Appendix~\ref{app:status}, row 1.1.}

\begin{remark}[Further consequences]\label{rem:post-fixed-cons}
(i) \emph{Sensitivity and specificity.} By Lemma~\ref{lem:summaries} and~\eqref{eq:post-exp}, inclusion probabilities of
true predictors are close to one and those of null predictors close to zero, at rate $\rn$, uniformly in $j$.
(ii) \emph{Without beta-min.} Under the other assumptions of Theorem~\ref{thm:post-fixed}, the posterior mass of models
that strictly contain $\Sz$ (over-fitted models) is $\OP(\rn)$, and every true predictor with
$\theta_j^{*2}\ge c\,\sigs^2\log n/(n\varphi_0^2)$, $c>c_0(\sz,A)$, has $1-\PIP_j=\OP(\rn)$. Specificity can fail
without beta-min (a numerical counterexample, reported with the proof, shows this).
(iii) \emph{Large models.} Models of size at least $3(A+1)(\sz+1)\log n$ carry posterior mass $\OP(n^{-1/2})$ with no
eigenvalue condition; this uses a bound on the coefficients instead.
\end{remark}

\status{(ii): part of the posterior concentration result, with the status of the in-probability form of
Theorem~\ref{thm:post-fixed}. (iii): the size-tail lemma, \textbf{proved, not yet independently checked}.}

Can the eigenvalue floor away from the true model, and the beta-min condition, be dropped?

\inwords{The conclusions of Theorem~\ref{thm:post-fixed} still hold if, instead of a floor on every small set of
predictors and a lower bound on each true coefficient, we ask only that the true predictors are not collinear among
themselves, and that every small model that misses a true predictor leaves unexplained signal of at least a constant times
$\sigs^2\log n$. The vanishing floor \as{NC} is kept.}

\begin{localassumption}[Floor on the true model and an identifiability margin]\label{asm:margin}
\begin{asmitems}
  \item\label{asm:margin-floor} $\lmin(\Rhat_{\Sz})\ge\varphi_0^2$ for a fixed $\varphi_0\in(0,1]$.
  \item\label{asm:margin-margin} Every model $S\in\Om$ that misses a true predictor and has at most
    $j_0:=\lceil(2A+1)\sz\rceil+1$ predictors outside $\Sz$ leaves unexplained signal
    $\norm{(I-P_S)\mu}^2\ge c_{\mathrm{id}}\,\sigs^2\log n$, for a fixed $c_{\mathrm{id}}>c_1(\sz,A)$, where $c_1(\sz,A)$ is an
    explicit constant that depends only on $(\sz,A)$.
\end{asmitems}
\rulesout{small wrong models that explain almost all of the mean signal.}
\end{localassumption}

Beyond the vanishing floor \as{NC}, Assumption~\ref{asm:margin} asks for no eigenvalue condition on sets other than the
true model.


\begin{theorem}[Posterior concentration under an identifiability margin]\label{thm:post-margin}
Let the design be fixed and assume \as{GL}, \as{DIM}, \as{SP}, \as{BC}, \as{NC} and Assumption~\ref{asm:margin}. Then
the conclusions~\eqref{eq:post-exp}--\eqref{eq:post-op} of Theorem~\ref{thm:post-fixed} hold, the expectation bound with
constants depending only on $(\sz,A,\sigs^2,\bar\theta,\varphi_0,g_0,c_{\mathrm{id}})$.
\end{theorem}

\status{\textbf{Proved and checked, conditional on the same two supporting results} (the lower-bound proposition
and the size-tail lemma, not yet independently checked) as the expectation bound of Theorem~\ref{thm:post-fixed}. \bridgepending{} Verbatim: Appendix~\ref{app:status}, row 1.1.}

\begin{remark}[On the assumptions]\label{rem:post-fixed-asm}
(i) \emph{The constants.} $c_0(\sz,A)$ and $c_1(\sz,A)$ are explicit but not optimised (Appendix~\ref{app:constants}). Under \as{EF} and \as{BM}(a),
Assumption~\ref{asm:margin} holds with $c_{\mathrm{id}}=C_\beta^2\varphi_0^2(1-1/n)$. So Theorem~\ref{thm:post-margin}
contains Theorem~\ref{thm:post-fixed} when $C_\beta^2\varphi_0^2>c_1$; when $c_0<C_\beta^2\varphi_0^2\le c_1$,
the two theorems coexist: Theorem~\ref{thm:post-fixed} applies, and Theorem~\ref{thm:post-margin} is not shown to apply.
(ii) \emph{Coefficients.} \as{BC} is used only through the energy bound $\norm{\mu}^2\le n\sz^2\bar\theta^2$; any bound
$\norm{\mu}^2\le nM_\mu$ with a fixed $M_\mu$ suffices.
(iii) \emph{A weaker form of \as{NC}.} \as{NC} can be weakened to a condition on the Occam factors of over-fitted
models, the \emph{superset Occam-sum condition} (stated in full with the proof); \as{NC} implies it for all large $n$. A
readable sufficient form is determinant growth: for a constant $c_D>0$ and every set $T$ with
$1\le\abs{T}\le3(A+1)(\sz+1)\log n$,
\[
  \frac{\det H_{\Sz\cup T}\,\det(G_{\Sz}+I)}{\det H_{\Sz}\,\det(G_{\Sz\cup T}+I)}\ \ge\ \{c_D(n+1)\}^{\abs{T}} .
\]
\end{remark}

\begin{remark}[What is not claimed]\label{rem:post-fixed-not}
(i) The floor on the true model in Assumption~\ref{asm:margin}\ref{asm:margin-floor} is sufficient and not known to be
necessary.
(ii) Theorem~\ref{thm:post-margin} keeps the vanishing floor \as{NC}, which still excludes exact linear dependence among
$O(\log n)$ predictors, duplicated predictors included.
(iii) No two-sided rate on every design is claimed. On every design the posterior miss is at least of order $n^{-1/2}$
(Theorem~\ref{thm:sharp}(a)); the extra $\sqrt{\log n}$ is attained on the designs of Theorem~\ref{thm:sharp}(b), and on
other designs it is not known.
\end{remark}

\begin{remark}[At moderate $n$]\label{rem:moderate}
The theorems describe the limit $n\to\infty$. The posterior miss is carried mainly by the $p-\sz$ models
$\Sz\cup\{j\}$ that add one null predictor to the true model. Each is individually light: its posterior odds relative to
$\Sz$ are of order $n^{-1/2}/p$, times a random factor driven by a $\chi^2_1$ variable. But they are numerous, and together
they can weigh as much as $\Sz$ itself. So at moderate $n$ and large $p$ the posterior is spread over $\Sz$ and many
models with one extra null predictor.
\end{remark}

\begin{remark}[Comparison with \citet{SongLiang2015}]\label{rem:SL-fixed}
Theorems~\ref{thm:post-fixed} and~\ref{thm:post-margin} give selection consistency for the full model, under the prior
as the method sets it and with no tuning: an explicit rate in expectation that is uniform over designs, and exact
recovery by the highest posterior model and the median probability model. Theorem~3 of Song and Liang gives sure
screening only: its sensitivity tends to one, and its specificity is not controlled; their consistency result for the
maximum a posteriori model is in their supplement. Theorem~\ref{thm:post-margin} replaces the floor \as{EF} on sets of
bounded size away from the true model, and the beta-min condition, by an identifiability margin on the residual signal,
while keeping the vanishing floor \as{NC}. This is in the style of \citet{NarisettyHe2014}; compare the population-level
identifiability condition of Song and Liang, their Condition~3, used for their Theorem~3. Our results are narrower in
dimension: we allow $p\le n^A$, while Song and Liang allow $P_n<e^{Cn^\alpha}$ with $\alpha<1$, through their
Corollary~1 and their condition~(10), which their Theorem~3 inherits.
\end{remark}

\subsection{Random Gaussian designs}\label{sec:random}

When the rows of the design are drawn from a Gaussian law, the sample conditions of Section~\ref{sec:fixed} hold with
probability tending to one under a condition on the population correlation matrix alone.

\inwords{If the rows are Gaussian, the true coefficients are fixed and nonzero, and no set of $O(\log n)$ predictors is
collinear in the population, then the posterior concentrates on the true model at the rate $\rn$, as for fixed designs,
now under the joint law of the design and the noise.}

\begin{assumption}{GR}{Gaussian rows}
The rows of $\Xr$ are i.i.d.\ $\Normal_p(m_n,\Sig)$, where $\Sig$ is a correlation matrix (unit diagonal).
\rulesout{dependent rows and non-Gaussian predictors.}
\end{assumption}

\begin{assumption}{FSb}{Fixed signals, raw scale}
$b_{\mathrm{lo}}\le\abs{\beta^*_j}\le b_{\mathrm{hi}}$ for every $j\in\Sz$, with fixed $0<b_{\mathrm{lo}}\le b_{\mathrm{hi}}<\infty$.
\rulesout{true coefficients that shrink to zero or blow up as $n$ grows.}
\end{assumption}

\begin{assumption}{EFS}{Population eigenvalue floor}
$\lmin(\Sigma_{n,S})\ge\varphi^2$ for every $S$ with $\abs{S}\le5(A+1)(\sz+1)\log n$, for a fixed $\varphi\in(0,1]$ (for
instance, $\lmin(\Sig)\ge\varphi^2$).
\rulesout{population collinearity within sets of $O(\log n)$ predictors.}
\end{assumption}

\begin{theorem}[Posterior concentration, random Gaussian design]\label{thm:post-random}
Assume \as{GL}, \as{DIM}, \as{SP}, \as{GR}, \as{FSb} and \as{EFS}. Then, under the joint law of the design and the noise,
\begin{gather}
  1-\pis{\Sz}=\OP\Big(\sqrt{\frac{\log n}{n}}\Big),\qquad \P\big(\HPM=\MPM=\Sz\big)\to1,\label{eq:post-random}\\
  \max_{j\in\set{p}}\big|\PIP_j-\ind{j\in\Sz}\big|=\OP\Big(\sqrt{\frac{\log n}{n}}\Big).\nonumber
\end{gather}
\end{theorem}

\status{The fixed-design result it rests on (Theorem~\ref{thm:post-fixed}, in-probability form): \textbf{proved and independently verified} (two independent verification passes: one by a reviewer of the same model
family as the author, one by a reviewer of a different model family; all inputs of the result were discharged on
2026-11-16). The
transfer from population to sample conditions: \textbf{proved and independently checked}. \bridgepending{} Verbatim: Appendix~\ref{app:status}, row 1.2.}

\begin{remark}[On the assumptions]\label{rem:post-random}
(i) No upper bound on the eigenvalues of $\Sig$ and no beta-min condition beyond \as{FSb} are needed, and $p\gg n$ is
allowed. The underlying result allows the floor in \as{EFS} to vanish slowly, as long as it is much larger than
$(\log n)^{3/2}/\sqrt n$; the statement uses a constant floor for simplicity.
(ii) Gaussian rows are used only to show that sample correlations of small sets of predictors are close to their
population values; sub-Gaussian rows are expected to work but are not claimed. Without the lower bound $b_{\mathrm{lo}}$
in \as{FSb}, the statements of Remark~\ref{rem:post-fixed-cons}(ii) still hold.
\end{remark}

\begin{remark}[Comparison with \citet{SongLiang2015}]\label{rem:SL-random}
This is the random-design setting of Song and Liang; here $p\gg n$ is allowed.
\end{remark}

\subsection{The rate cannot be improved}\label{sec:sharp}

Is $\rn$ the right rate? This subsection shows that no design allows a rate faster than $n^{-1/2}$, and that on two
families of designs the extra factor $\sqrt{\log n}$ is really there.

\inwords{(a) Whatever the design, the posterior cannot put more than about $1-c/\sqrt n$ of its mass on any one
model with at most $\min(K-1,p/2)$ predictors, because the prior and the Occam factor keep every one-predictor extension of it alive. (b) On designs whose null
predictors are orthogonal, or independent Gaussian, the posterior miss is on average at least a constant times
$(\sz+1)\rn$, for every number of predictors. So the rate of Theorem~\ref{thm:post-fixed} cannot be improved under its
assumptions.}

The lower bounds in part (b) hold on the following \emph{witness} designs. They are not conditions we need for
consistency; they describe where the rate is attained.

\begin{localassumption}[Witness designs]\label{asm:witness}
\begin{asmitems}
  \item\label{asm:witness-energy} The true model size $\sz$ is fixed, and $\norm{\mu}^2\le nM_\mu$ and
    $\norm{\tstar}^2\le nM_\theta$ for fixed $M_\mu$, $M_\theta$.
  \item\label{asm:witness-orth} \emph{Orthogonal null predictors.} The columns outside $\Sz$ are mutually orthogonal
    and orthogonal to the columns in $\Sz$.
  \item\label{asm:witness-gauss} \emph{Gaussian null predictors.} The raw columns outside $\Sz$ have i.i.d.\ $\Normal(0,1)$
    entries, independent of $\Xr_{\Sz}$ and of $\varepsilon$, and $p\le n^A$.
\end{asmitems}
\end{localassumption}

\begin{theorem}[The rate cannot be improved]\label{thm:sharp}
\textup{(a)} \emph{Every design.} For every $n\ge3$, every standardised design with $p\ge2$, every response and every
model $S$ with $\abs{S}\le\min(K-1,p/2)$,
\begin{equation}\label{eq:sharp-a}
  \pis{S}\ \le\ \Big(1+\frac{\abs{S}+1}{3\sqrt n}\Big)^{-1}.
\end{equation}
\textup{(b)} \emph{Witness designs.} Assume \as{GL}, Assumption~\ref{asm:witness}\ref{asm:witness-energy}, and either
Assumption~\ref{asm:witness}\ref{asm:witness-orth} or~\ref{asm:witness-gauss}. Then, for every $p>\sz$, bounded or
growing,
\begin{equation}\label{eq:sharp-b}
  \E\big[1-\pis{\Sz}\big]\ \ge\ (1-o(1))\,(\sz+1)\,\frac{2(p-\sz)}{2p-1-\sz}\,
  \sqrt{\frac{\log\!\big(p\sqrt n/(\sz+1)\big)}{\pi n}},
\end{equation}
and, if moreover $p\to\infty$,
\begin{equation}\label{eq:sharp-c}
  1-\pis{\Sz}\ \ge\ (1-\oP(1))\,(\sz+1)\sqrt{\frac{\log p}{\pi n}}.
\end{equation}
\end{theorem}

\status{\textbf{Proved and independently checked}. \bridgepending{} Verbatim: Appendix~\ref{app:status}, row 1.3.}

In~\eqref{eq:sharp-b} and~\eqref{eq:sharp-c}, $\pi$ is the number $3.14\ldots$, not the posterior.
Display~\eqref{eq:sharp-a} says that no model of that size, the true one included, can carry posterior mass closer to one
than about $(\abs{S}+1)/(3\sqrt n)$. Display~\eqref{eq:sharp-b} says that, on the witness designs, the expected posterior
miss is at least $c\cdot(\sz+1)\cdot\rn$, where $c>0$ depends only on $\sz$, and does not depend on $\sz$ when
$\sz\le p/2$.

\begin{remark}[Further consequences]\label{rem:sharp}
(i) \emph{Size of the true model.} By~(a), for true models with $\sz\le\min(K-1,p/2)$, no design gives a rate faster than
$n^{-1/2}$; $\pis{\Sz}\to1$ requires $\sz=o(\sqrt n)$; and the rate $\rn$ requires $\sz=O(\sqrt{\log n})$. (Without the
size restriction, (a) is false: for example $S=\set{p}$ with $p=2$.)
(ii) \emph{Where the $\sqrt{\log n}$ appears.} On the witness designs it is present in expectation for every $p$, bounded
$p$ included (with $p=n^a$ the constant is $\sqrt{(a+\frac12)/\pi}$), and in probability whenever $\log p\asymp\log n$. In
expectation at bounded $p$ it comes from rare large gains of a single extra predictor, not from the number of
competitors.
(iii) \emph{Exactness.} The rate $\rn$ of Theorem~\ref{thm:post-fixed} is exact in two named cases. In expectation: on
fully orthogonal designs (all columns mutually orthogonal, any $p$ with $\sz<p\le n-1$, bounded or growing) with fixed
nonzero coefficients; these designs satisfy the assumptions of Theorem~\ref{thm:post-fixed} and of~\eqref{eq:sharp-b},
so $\E[1-\pis{\Sz}]$ is of exact order $\rn$, with constants depending only on $\sz$, the coefficients and the noise
variance (conditional on Theorem~\ref{thm:post-fixed}). In probability: on fully orthogonal designs with
$p=\min(n-1,\lfloor n^a\rfloor)$, and on raw designs with i.i.d.\ $\Normal(0,1)$ entries and $p=\lfloor n^a\rfloor$, both
with fixed nonzero coefficients (conditional on Theorems~\ref{thm:post-fixed} and~\ref{thm:post-random}).
\end{remark}

\begin{remark}[Further consequences: growing true models]\label{rem:sharp-growing}
(i) Keep the energy bounds of Assumption~\ref{asm:witness}\ref{asm:witness-energy}, let all columns outside $\Sz$ be null
columns as in Assumption~\ref{asm:witness}\ref{asm:witness-orth}, or as in Assumption~\ref{asm:witness}\ref{asm:witness-gauss} with the columns in $\Sz$ and the coefficients
fixed (nonrandom), and let $\sz\le p/2$ grow with $n$. Then a posterior mean error $\E[1-\pis{\Sz}]$ of order $\rn$
forces $\sz$ to stay bounded, for every $p$.
(ii) \emph{Critical scale} (orthogonal null predictors only). Keep the energy bounds and Assumption~\ref{asm:witness}\ref{asm:witness-orth},
let $p\to\infty$, and let $\sz$ grow with $\sz/p\to0$ and $(\sz+1)\sqrt{\log p/(\pi n)}\to\lambda\in(0,\infty)$. Then
$\pis{\Sz}\le e^{-\lambda}+\oP(1)$.
\end{remark}

\begin{remark}[What is not claimed]\label{rem:sharp-not}
(i) On a general design only the order $n^{-1/2}$ of part~(a) is proved as a lower bound; a matching lower bound
$\E[1-\pis{\Sz}]\gtrsim\rn$ under a floor on the sets $\Sz\cup\{j\}$ is open.
(ii) When $\log p=o(\log n)$, nothing is claimed in probability: \eqref{eq:sharp-c} gives only the order
$\sqrt{\log p/n}$, and whether the $\sqrt{\log n}$ is present in probability there is open.
(iii) On raw Gaussian designs the matching upper bound in expectation is open (Theorem~\ref{thm:post-random} is an
in-probability statement), and whether the constant in~\eqref{eq:sharp-b} is attained is not proved.
\end{remark}

\subsection{Omitted variables}\label{sec:omitted}

In practice some true predictors may be missing from the data. Here the method is run on a set of candidate predictors
only, and we ask what the posterior then selects.

\inwords{If the method sees only some of the predictors, and the response given those predictors is still Gaussian and
linear in them, the posterior selects the predictors that the best linear predictor built from them uses, with error of
order $\rn$. The residual variance may be fixed, may vanish, or may grow with $n$.}

In this subsection the method is run on the response $\yr$ and a set $\mathcal C$ of candidate predictors, with rows
$x_{\mathcal C,1},\dots,x_{\mathcal C,n}$; $p$ is the number of candidates, and $\Sigma$ and $R$ are the covariance and the
correlation matrix of a candidate row.

\begin{localassumption}[Candidate-only fitting]\label{asm:omitted}
\begin{asmitems}
  \item\label{asm:omitted-model} \emph{Model.} The candidate rows are i.i.d.\ $\Normal_p(m,\Sigma)$, and, given the
    candidate design, the responses are independent with
    \[
      \yr_i\mid x_{\mathcal C}\ \sim\ \Normal\big(\alpha'+x_{\mathcal C,i}^\top\beta^{\mathcal C},\ \sigma^2_{\mathcal C}\big),
      \qquad \sigma^2_{\mathcal C}>0 .
    \]
    We write $\Sz^{\mathcal C}:=\supp\beta^{\mathcal C}$ and $s:=\abs{\Sz^{\mathcal C}}$.
  \item\label{asm:omitted-floor} \emph{Floor.} $\lmin(R_{TT})\ge c_{\mathcal C}$ for every $T\subseteq\mathcal C$ with
    $\abs{T}\le3(A+1)(s+1)\log n$, for a fixed $c_{\mathcal C}>0$.
  \item\label{asm:omitted-signal} \emph{Signal relative to the residual variance.}
    $\beta^{\mathcal C\top}\Sigma\beta^{\mathcal C}\le M_{\mathcal C}\,\sigma^2_{\mathcal C}$, and
    \[
      \frac{n\min_{j\in\Sz^{\mathcal C}}\Sigma_{jj}(\beta^{\mathcal C}_j)^2}{\sigma^2_{\mathcal C}\log n}\to\infty .
    \]
\end{asmitems}
\rulesout{(a) a response whose law given the candidates is not Gaussian with a mean linear in them and a constant
variance; (b) population collinearity among $O(\log n)$ candidates; (c) candidate signal that is huge compared with the
residual variance, and candidate coefficients too small to detect at the scale $\rn$.}
\end{localassumption}

In words, $\beta^{\mathcal C}$ are the coefficients of the best linear predictor of the response from the candidates,
and $\sigma^2_{\mathcal C}$ is the variance left over, which includes the contribution of the omitted predictors.

\begin{theorem}[Omitted variables]\label{thm:omitted}
Assume \as{DIM} (with $p$ the number of candidates), Assumption~\ref{asm:omitted}, and that $s\ge1$ is fixed. No
condition is placed on how $m$, $\Sigma$, $\beta^{\mathcal C}$ or $\sigma^2_{\mathcal C}$ vary with $n$. Then
\begin{equation}\label{eq:omitted}
  \E\big[1-\pi(\Sz^{\mathcal C}\mid y)\big]=O\Big(\sqrt{\frac{\log n}{n}}\Big),\qquad
  \P\big(\HPM\neq\Sz^{\mathcal C}\big)=O\Big(\sqrt{\frac{\log n}{n}}\Big).
\end{equation}
\end{theorem}

\status{\textbf{Proved and checked, conditional on} the expectation bound of Theorem~\ref{thm:post-fixed} (hence on its
two supporting results, the lower-bound proposition and the size-tail lemma, not yet independently checked) and on two
lemmas on Gaussian designs that are verified and passed a council review. \bridgepending{} Verbatim: Appendix~\ref{app:status}, row 1.4.}

\begin{example}[Omitted Gaussian predictors]\label{ex:omitted}
Suppose the rows $(x_{\mathcal C,i},x_{\mathcal M,i})$ of candidates and omitted predictors are i.i.d.\ Gaussian with
$\Sigma_{\mathcal C\mathcal C}$ invertible, and
$\yr_i=\alpha+x_{\mathcal C,i}^\top\beta_{\mathcal C}+x_{\mathcal M,i}^\top\beta_{\mathcal M}+\varepsilon_i$ with
$\varepsilon\sim\Normal_n(0,\sigma^2I)$ independent of the rows. Then the full model is misspecified for the
candidate-only fit, and Assumption~\ref{asm:omitted}\ref{asm:omitted-model} holds with
\[
  \beta^{\mathcal C}=\beta_{\mathcal C}+\Sigma_{\mathcal C\mathcal C}^{-1}\Sigma_{\mathcal C\mathcal M}\beta_{\mathcal M},
  \qquad
  \sigma^2_{\mathcal C}=\sigma^2+\beta_{\mathcal M}^\top\big(\Sigma_{\mathcal M\mathcal M}
  -\Sigma_{\mathcal M\mathcal C}\Sigma_{\mathcal C\mathcal C}^{-1}\Sigma_{\mathcal C\mathcal M}\big)\beta_{\mathcal M}.
\]
The first formula says that each candidate absorbs the part of the omitted signal it can predict; the second, that the
rest of the omitted signal is added to the noise. If $\Sigma_{\mathcal C\mathcal C}$ is singular, the same holds with its
Moore--Penrose inverse. The same holds when the omitted predictors are a linear function of the candidates plus Gaussian
noise whose rows are i.i.d.\ and independent of the candidates.
\end{example}

\begin{remark}[Further consequences]\label{rem:omitted-cons}
(i) \emph{What the posterior targets.} The law in Assumption~\ref{asm:omitted}\ref{asm:omitted-model} is the conditional
law of $\yr$ given the candidates, hence the Kullback--Leibler projection of the truth onto the candidate linear models:
the posterior concentrates on the support of the best candidate predictor.
(ii) \emph{Contraction.} Under the same assumptions the posterior's fitted mean and variance also contract around this
law, in the in-sample average Hellinger and Kullback--Leibler distances, at rate $M_n/\sqrt n$ in Hellinger and $M_n/n$
in Kullback--Leibler, for any $M_n\to\infty$. Here the mean of the law is centred over the sample, because the fitted mean
carries no intercept.
For a fixed design the contraction needs only that the posterior mass of $\Sz^{\mathcal C}$ tends to one in probability
(with the energy bound and a floor on $\Sz^{\mathcal C}$); it holds uniformly over any class of fixed-design
configurations with common constants and bounded support size on which this convergence is uniform.
(iii) \emph{Fixed design.} A fixed-design version of the theorem, uniform over the residual variance, also holds.
\end{remark}

\status{Remark~\ref{rem:omitted-cons}(ii): \textbf{proved and checked}, conditional on the posterior contraction lemma for
fixed designs (itself proved and checked) and, for the random-design clause, on the same inputs as
Theorem~\ref{thm:omitted}. (iii): \textbf{proved and checked, conditional on} the expectation bound of
Theorem~\ref{thm:post-fixed}. Verbatim: Appendix~\ref{app:status}, row 1.4.}

\begin{remark}[On the assumptions]\label{rem:omitted-asm}
(i) \emph{A weaker sufficient condition.} The proof uses only that, given the candidates, the coordinates of the residual
vector $\yr-\alpha'\one-x_{\mathcal C}\beta^{\mathcal C}$ in an orthonormal basis of the vectors orthogonal to $\one$ are
i.i.d.\ $\Normal(0,\sigma^2_{\mathcal C})$; nothing is required of the residual's sample mean. The Gaussian law of the
candidate rows is used only to transfer the population floor and signal conditions to the sample.
(ii) \emph{Why relative to $\sigma^2_{\mathcal C}$.} Everything is measured relative to $\sigma^2_{\mathcal C}$ because
the method is invariant to shifts and rescaling of $y$. The energy bound in
Assumption~\ref{asm:omitted}\ref{asm:omitted-signal} cannot simply be dropped (numerical evidence); the beta-min order is
that of Theorem~\ref{thm:post-fixed}.
(iii) \emph{Relation to Theorem~\ref{thm:post-random}.} Theorem~\ref{thm:post-random} does not give
Theorem~\ref{thm:omitted}: it fixes the noise variance and has no omitted predictors. Whether the hypotheses of
Theorem~\ref{thm:omitted} cover the setting of Theorem~\ref{thm:post-random} without omitted predictors is open. For now the two
theorems rest on different arguments, each with its own conditions.
\end{remark}

\begin{remark}[Comparison with \citet{SongLiang2015}]\label{rem:SL-omitted}
Song--Liang (2015, Theorem~4, eqs.~(17)--(18); Corollary~2) fit a fixed candidate subset that misses true predictors, a
model they call misspecified, with candidates having finite second moments only (their Condition~4): the posterior is
Hellinger-consistent at the Kullback--Leibler projection, for each fixed $\varepsilon$, almost surely and without a rate,
and, if the candidates have no exact collinearity, the posterior probability that a sampled model contains the
projection's support tends to one in probability (Corollary~2). Here, assuming instead that, given the candidates, the
response is Gaussian with a mean linear in them and a constant variance (as when the omitted predictors are jointly
Gaussian with the candidates), we prove exact selection of the best candidate linear predictor's support (Song--Liang's
$\tilde s$ for centred rows) at rate $\rn$, with up to $n^A$ candidates and no condition on how $\Sigma$, $\beta$,
$\sigma^2$ vary with $n$; the coded method is exactly invariant to shifts and rescaling of $y$.
\end{remark}
